\subsection{One transform, two kinds of data}
\source{This perspective is guided by Persiflage \cite{persiflage}.
The underlying transform and orthogonal-polynomial identities are
Pr\'evost's \cite[Section 2 and Theorems 1--2]{prevost}.}

It is helpful to see where the zeta value enters.
Change variables \(t=y^2\) and write
\[
 d\nu(t)=\frac{w(\sqrt t)}{2\sqrt t}\,dt,\qquad
 \Phi(z)=\int_0^\infty\frac{d\nu(t)}{z+t}\quad(z>0).
\]
The moments \(m_j=\int t^j\,d\nu(t)=\mu_X(t^j)\) are rational.
On the other hand, the pole formula gives
\[
 \Phi(a^2)=a^4\bigl(\xi-H_a^{(5)}\bigr)-\frac14+\frac1{2a}
 \qquad(a\in\Z_{>0}).
\]
Thus rational approximation to a transform value gives rational
approximation to \(\xi\).

The moments also describe \(\Phi\) for large \(z\).
The finite geometric identity gives
\[
 \Phi(z)=\sum_{j=0}^{m-1}\frac{(-1)^jm_j}{z^{j+1}}
  +\frac{(-1)^m}{z^m}\int_0^\infty\frac{t^m}{z+t}\,d\nu(t).
\]
The last term has absolute value at most \(m_m z^{-m-1}\).
This is an asymptotic expansion for each fixed \(m\); we have not asserted
convergence of the infinite moment series.

At \(X=\xi\), the two clauses defining \(\mu_X\) describe samples and
asymptotic coefficients of one transform. This is the bridge to
Pr\'evost's Pad\'e approximation, emphasized by Persiflage.
Our proof uses a rationally modified measure and its determinant;
no quantitative Pad\'e estimate is needed. The question is how much
real smallness survives when the determinant's coefficients are made
integral.
